\begin{textsample}{NEG Dim 3 – vem\_ed\_11 – Score 4.00 – vem\_ed\_11/t048.txt}  \label{ex:f3_neg_046}
PCIs com Índia valorizam diversidade cultural Em 2021, a equipe dos PCIs / Cesu iniciou colaborações com o Symbiosis Centre for Management Studies, localizado em Pune, oeste da Índia.
Os principais desafios contornados foram o fuso horário, com 8 horas e meia de diferença, e o tamanho das turmas indianas – existem salas com mais de 100 alunos.
Lark foi a principal ferramenta de comunicação utilizada, para realização dos encontros síncronos e compartilhamento de informações.
Além disso, utilizou-se bastante o WhatsApp para interação entre os participantes.
Dois projetos foram realizados no segundo semestre de 2021: um sobre satisfação no trabalho, nas Fatecs Americana (Ricardo Bertoni Pompeu), Sumaré (Danilo Sergio Sorroce e Rafaeli Cardozo Modolo Begalli) e, na instituição indiana, a professora Nehajoan Panackal.
Nesse projeto participaram cerca de 50 alunos de cada país.
Situações-problema sobre satisfação no trabalho eram apresentadas em vídeo e os grupos propunham soluções. “Fomos surpreendidos pela qualidade do trabalho que encontramos com esse novo parceiro”, comenta Ana Carolina Freschi, da equipe dos PCIs, que acompanhou o projeto. “Tudo correu muito bem, Neha é uma professora muito engajada; os alunos tiveram uma interação muito grande.
Os indianos foram extremamente generosos”, comemora Pompeu.
O outro projeto, sobre impacto da Covid-19 em padrões de investimentos pessoais, foi conduzido pelos professores Ricardo Trovão (Fatec Itaquaquecetuba) e Shreya Virani (Symbiosis).
Os PCIs foram muito bem-sucedidos e bem avaliados, tanto pelos docentes brasileiros quanto pelas professoras indianas.
Portanto, as colaborações continuam em 2022 e se estendem para outras áreas como marketing.
Para além dos conteúdos técnicos, houve o enriquecimento cultural dos estudantes.
Durante a realização do PCI, em 4 de novembro, ocorreu o Diwali – festival hindu que celebra as divindades femininas Kali, Lakshmi e Sarasvati, culminando com a vitória das luzes sobre as trevas.
Os estudantes indianos se preocuparam em preparar uma apresentação sobre o Diwali para os brasileiros, contando sobre comidas, vestimentas, músicas. “Em poucas semanas, brasileiros e indianos conseguiram resolver problemas, interagir, trocar informações.
Foi um sucesso”, afirma Freschi.

% matched lemmas: 
\end{textsample}
